%-*-latex-*-
%-*-latex-*-
\newcommand\COURSE{ciss350}
\newcommand\ASSESSMENT{q10704}
\newcommand\ASSESSMENTTYPE{Quiz}
\newcommand\POINTS{\textwhite{xxx/xxx}}

\input{myquizpreamble}
\input{yliow}
\input{\COURSE}
%-*-latex-*-

%-*-latex-*-

%-*-latex-*-
\renewcommand\TITLE{\ASSESSMENTTYPE \ \ASSESSMENT}

\renewcommand\EMAIL{}
\renewcommand\AUTHOR{eagraham1@cougars.ccis.edu}

\textwidth=6in
\begin{document}
\topmatterthree

The following is the node for the binary search tree.
\begin{Verbatim}[frame=single,fontsize=\footnotesize]
class BSTNode
{
public:
    int key_;
    BSTNode * parent_;
    BSTNode * left_, * right_;
};
\end{Verbatim}

%------------------------------------------------------------------------------
\nextq
Complete the function below to return the address
of the BST node with the given target key value.
If it is not found, return \verb!NULL!.
Calling \verb!bst_find(p, 42)! will search for the node with
key value 42 among all the children of the node with address \verb!p!.
(This is just our usual BST find.)
\\
\ANSWER
\begin{answercode}
BSTNode * bst_find(BSTNode * p, int target)
{
  if(p == NULL)
    return NULL;

  if(p->key_ == target)
    return p;

  else
  {
    if(p->key_ < target)
      return bst_find(p->right_, target);
    else
      return bst_find(p->left_, target);
  }
}
\end{answercode}

%------------------------------------------------------------------------------
\nextq
Complete the function below to return the address
of the \textit{parent} node of the BST node with the given target key value.
If it is not found, return \verb!NULL!.
(There are two cases where \verb!NULL! is returned:
when the target key is at the root of the tree and when the target key
is not in the tree.)
\\
\ANSWER
\begin{answercode}
BSTNode * bst_find_parent(BSTNode * p, int target)
{
  if(p == NULL)
    return NULL;

  if(p->key_ == target)
    return p->parent_;

  else
  {
    if(p->key_ < target)
      return bst_find_parent(p->right_, target);
    else
      return bst_find_parent(p->left_, target);
  }
}
\end{answercode}

%------------------------------------------------------------------------------
\nextq
Complete the following function which is similar to the BST find.
If the target key value is found, the address of this target node is
returned and \verb!found! (pass by reference) is set to \verb!true!.
If the target key value is not found, the address of the last node
reached is returned and \verb!found! is set to \verb!false!.
For instance when you perform this version of find
on the following tree starting at \verb!5!
and searching for target value \verb!7!
\begin{Verbatim}[frame=single,fontsize=\footnotesize]
      (5)
     /   \
  (3)     (8)
         /   \
      (6)     (9)
\end{Verbatim}
the address returned is the address of the node with value \verb!6!
and \verb!found! is set to \verb!false!.
(This function is helpful: after failing to find a target key value
and you want to insert this new key,
it tells you where is the parent node for the new node.)
\\
\ANSWER
\begin{answercode}
BSTNode * bst_find_up_to(BSTNode * p, int target, bool & found)
{
  if(p == NULL)
  {
    found = false;
    return NULL;
  }

  if(p->key_ == target)
  {
    found = true;
    return p;
  }

  if(p->key_ < target)
  {
    if(p->left_ == NULL)
    {
      found = false;
      return p;
    }
    return bst_find_up_to(p->right_, target, found);
  }
  
  else
  {
    if(p->right == NULL)
    {
      found = false;
      return p;
    }
    return bst_find_parent(p->left_, target, found);
  }
}
\end{answercode}

%------------------------------------------------------------------------------
%-*-latex-*-
\newpage
\input{instructions}
\end{document}

